\documentclass[10pt]{article}
\usepackage[T1]{fontenc}
\usepackage[utf8]{inputenc}
\usepackage[french]{babel}
\usepackage{amsmath,amssymb,mathrsfs,mathpazo}
\usepackage{geometry}
\geometry{a4paper,margin=2cm}
\usepackage{graphicx,array,tabularx,tabularray,colortbl,multirow,multicol,booktabs,enumitem,pifont,ulem,fancybox,diagbox,cancel}
\usepackage{pstricks,pstricks-add,pst-plot,pst-tree,pst-node,pst-text,pst-func,pst-math,pst-bspline,pst-3dplot,pst-eucl,pst-solides3d}
\let\pscalculate\undefined
\usepackage{pst-bezier}
\usepackage{tikz,pgfplots,tkz-tab}
\usetikzlibrary{arrows,arrows.meta,patterns,calc,positioning,trees,decorations.markings,shapes.geometric,shapes.misc,shapes.arrows,intersections,decorations.pathmorphing,decorations.pathreplacing,backgrounds,shadows,babel}
\pgfplotsset{compat=1.18}
\usepgfplotslibrary{fillbetween}
\usepackage[np]{numprint}
\usepackage{listings,makecell,pas-tableur,twoopt}
\let\barre\undefined
\usepackage[pstricks]{tablvar}
\definecolor{darkgreen}{RGB}{0,100,0}
\providecommand{\red}{\color{red}}
\providecommand{\blue}{\color{blue}}
\providecommand{\green}{\color{green}}
\providecommand{\black}{\color{black}}
\providecommand{\white}{\color{white}}
\providecommand{\gray}{\color{gray}}
\providecommand{\cyan}{\color{cyan}}
\providecommand{\magenta}{\color{magenta}}
\providecommand{\yellow}{\color{yellow}}
\providecommand{\violet}{\color{violet}}
\providecommand{\orange}{\color{orange}}
\providecommand{\euro}{\texteuro}
\providecommand{\EUR}{\texteuro}
\providecommand{\ieme}{\textsuperscript{e}}
\providecommand{\degres}{\ensuremath{^\circ}}
\providecommand{\degree}{\ensuremath{^\circ}}
\providecommand{\vv}[1]{\overrightarrow{#1}}
\providecommand{\Coord}[1]{#1}
\providecommand{\decosix}{}
\providecommand{\starredbullet}{\ensuremath{\bullet}}
\pagestyle{empty}
\setlength{\parindent}{0pt}
\usepackage{adjustbox}\DeclareUnicodeCharacter{2236}{:}\DeclareUnicodeCharacter{2190}{\ensuremath{\leftarrow}}
\providecommand{\R}{\mathbb{R}}
\providecommand{\N}{\mathbb{N}}
\providecommand{\Z}{\mathbb{Z}}
\providecommand{\Q}{\mathbb{Q}}
\providecommand{\D}{\mathbb{D}}
\providecommand{\C}{\mathbb{C}}
\providecommand{\e}{{\mathrm{e}}}
\providecommand{\vect}[1]{\overrightarrow{#1}}
\providecommand{\vectt}[1]{\overrightarrow{\mathrm{#1}}}
\providecommand{\vvt}[1]{\overrightarrow{\mathrm{#1}}}
\providecommand{\np}[2][]{{#2~#1}}
\providecommand{\vv}[1]{{\overrightarrow{#1}}}
\renewcommand{\d}{{\mathrm{d}}}
\providecommand{\degres}{{^\circ}}
\providecommand{\degre}{{^\circ}}
\providecommand{\blue}{\color{blue}}
\providecommand{\red}{\color{red}}
\providecommand{\enskip}{\quad}
\providecommand{\llbracket}{\lbrack\!\lbrack}
\providecommand{\rrbracket}{\rbrack\!\rbrack}
\providecommand{\euro}{{\text{€}}}
\providecommand{\no}{{\text{n°}}}
\providecommand{\interval}[2]{{[#1;#2]}}
\providecommand{\intervFO}[2]{{[#1;#2[}}
\providecommand{\dotfill}{\ldots}

\providecommand{\ds}{\displaystyle}
\providecommand{\textding}[1]{#1}
\providecommand{\Oij}{\(\mathrm{O};\vec{\imath},\vec{\jmath}\)}
\providecommand{\Oijk}{\(\mathrm{O};\vec{\imath},\vec{\jmath},\vec{k}\)}
\providecommand{\Ouv}{\(\mathrm{O};\vec{u},\vec{v}\)}
\providecommand{\intervFF}[2]{\ensuremath{[#1;#2]}}
\renewcommand{\euro}{\ifmmode\text{\texteuro}\else\texteuro\fi}
\DeclareUnicodeCharacter{00B0}{\ifmmode^{\circ}\else\textdegree\fi}
\DeclareUnicodeCharacter{20AC}{\ifmmode\text{\texteuro}\else\texteuro\fi}
% Énoncés : Éducation nationale. Transcription et corrigés : APMEP, Denis Vergès et contributeurs. Découpage et classement : CoopMaths, Sébastien Lozano, Éric Elter et contributeurs. Adaptation de la présentation : Maths Explorer.
% CC BY-SA : https://coopmaths.fr/www/about/
% https://coopmaths.fr/alea/static/eam/2026/tex/eam_specifique_2026_06_antilles_2_cor.tex
\begin{document}
\noindent\textit{APMEP / CoopMaths. CC BY-SA. Présentation adaptée par Maths Explorer.}\par\medskip
On compte le nombre d'arbres de deux forêts où de nouveaux arbres sont plantés chaque année.

\subsubsection*{Partie A\hspace{10mm}Première forêt}

Dans cette forêt, au 1\up{er} janvier 2010, on dénombrait \np{1200}~arbres.

On suppose que le nombre d'arbres augmente de \np{100} par an.

On modélise cette évolution par une suite \mbox{$(u_n)$}, où $u_n$ représente le nombre d'arbres dans cette forêt au 1\up{er}~janvier de l'année \mbox{$2010+n$} pour tout entier naturel $n$.

Ainsi, on a \mbox{$u_0=\np{1200}$}
\begin{enumerate}
	\item Calculer $u_1$ et $u_2$. Interpréter $u_2$ dans le contexte de l'exercice.


	\textcolor{red}{\textsc{Correction}}\\
	Le nombre d'arbres augmente de 100 par an, donc \mbox{$u_1=\np{1200}+100=\np{1300}$}  et \mbox{$u_2=\np{1300}+100=\np{1400}$}.

	Dans le contexte de l'exercice, cela signifie que la forêt comptera \np{1400}~arbres en 2012.
	\item Exprimer \mbox{$u_{n+1}$} en fonction de $u_n$ pour tout entier naturel $n$. En déduire la nature de la suite $(u_n)$ et préciser sa raison $r$.


	\textcolor{red}{\textsc{Correction}}\\
	Chaque année le nombre d'arbres augmente de 100, donc \mbox{$\forall n \in \N,~ u_{n+1}=u_n + 100$}.

	La suite $(u_n)$ est donc une suite arithmétique, de raison 100 et de premier terme \mbox{$u_0=\np{1200}$}.
		\item Exprimer $u_n$ en fonction de $n$.


		\textcolor{red}{\textsc{Correction}}\\
		L'expression explicite de $(u_n)$ est donc  \mbox{$\forall n \in \N,~ u_n=u_0 +r\times n= \np{1200} + 100~n$}
	\item D'après ce modèle, à partir de quelle année la forêt comptera-t-elle plus de \np{2950}~arbres ?

		\textcolor{red}{\textsc{Correction}}\\
		Comme le nombre d'arbres est toujours un multiple de 100, la forêt comptera plus de \np{2950}~arbres\dots lorsqu'elle en aura \np{3000}.

		\mbox{$u_n=\np{3000}\iff \np{1200}+100~n=\np{3000}\iff 100~n=\np{1800}$}, donc il faut que \mbox{$n=18$} : c'est à partir de 2028 que la forêt dépassera \np{2950}~arbres.
\end{enumerate}

\subsubsection*{Partie B\hspace{10mm}Seconde forêt}


Dans cette forêt, au 1\up{er} janvier 2010, on dénombrait \np{1000}~arbres.

On suppose que le nombre d'arbres augmente de 5\% par an.

On modélise cette évolution par une suite \mbox{$(v_n)$}, où $v_n$ représente le nombre d'arbres dans cette forêt au 1\up{er}~janvier de l'année \mbox{$2010+n$} pour tout entier naturel $n$.

Ainsi, on a \mbox{$v_0=\np{1000}$}

\begin{enumerate}
	\item Calculer le nombre d'arbres au 1\up{er}~janvier 2011.


		\textcolor{red}{\textsc{Correction}}\\
		Le nombre d'arbres augmente de 5\,\% par an, et 5\,\% de \np{1000} c'est 50.

		Donc au 1\up{er} janvier 2011, il y aura dans cette forêt \np{1050}~arbres.
	\item Déterminer la nature de la suite \mbox{$(v_n)$} et préciser sa raison $r$.

		\textcolor{red}{\textsc{Correction}}\\
		Chaque année, le nombre d'arbres augmente de 5\,\%, ce qui revient à multiplier la quantité de l'année précédente par \np{1,05}.

		$(v_n)$ est donc une suite géométrique, de raison \np{1,05} et de premier terme \np{1000}.
	\item D'après ce modèle, on estime qu'au 1\up{er}~janvier 2030, il y aura environ \np{2653}~arbres sur la parcelle. Quel terme de la suite a-t-on calculé? Quel a été le calcul effectué ?

		\textcolor{red}{\textsc{Correction}}\\
		L'année 2030 correspond à \mbox{$v_{20}$}, car \mbox{$2030=2010+20$.}


		L'expression explicite de la suite \mbox{$(v_n)$} est \mbox{$v_n=v_0\times q^n= \np{1000}\times \np{1.05}^n$}, donc on a effectué le calcul \mbox{$\np{1000}\times\np{1.05}^{20}$} pour trouver ce résultat.
\end{enumerate}
\subsubsection*{Partie C\hspace{10mm}Comparaison des évolutions}
Grâce à l'extrait de tableur ci-dessous, on visualise le nombre d'arbres arrondi à l'unité dans les deux forêts.
\begin{center}
	\begin{tblr}{|X[8mm]|*{9}{X[c]|}}\hline
		$n$&	0&	1&	2&	3&	\dots&	26&	27&	28&	\textbf{29}\\\hline
		$u_n$&	\np{1200}&	\np{1300}&	\np{1400}&	\np{1500}&	\dots&	\np{3800}&	\np{3900}&	\np{4000}&	\textbf{\np{4100}}\\\hline
		$v_n$&	\np{1000}&	\np{1050}&	\np{1103}&	\np{1158}&	\dots&	\np{3556}&	\np{3733}&	\np{3920}&\textbf{	\np{4116}}\\\hline
	\end{tblr}
	\end{center}
Interpréter dans le contexte de l'exercice les résultats de la dernière colonne.

	\textcolor{red}{\textsc{Correction}}\\
	Dans la dernière colonne, on observe que \mbox{$v_{29} > u_{29}$}, alors que pour tous les termes précédents, on  avait~\mbox{$u_n> v_n$}.

	On peut en déduire qu'à partir de 2039, la forêt n°2 aura un plus grand nombre d'arbres que la forêt~n°1.

\end{document}
